% Part: first-order-logic
% Chapter: sequent-calculus

\documentclass[../../../include/open-logic-chapter]{subfiles}

\begin{document}

\iftag{FOL}
      {\olchapter{fol}{seq}{क्रमवर्ती कलन}}
      {\olchapter{pl}{seq}{क्रमवर्ती कलन}}

\begin{editorial}
  या प्रकरणात अभिजात प्रथम-क्रम तर्कशास्त्रासाठी गेंटझेन यांचे मानक
  क्रमवर्ती कलन LK मांडले आहे. यात आणखी उदाहरणे आणि सराव समाविष्ट करणे
  उपयुक्त ठरेल. सिद्धता-पद्धती म्हणून क्रमवर्ती कलनाशी संबंधित मजकूर
  समाविष्ट करण्यासाठी किंवा वगळण्यासाठी ``prfLK'' हा टॅग वापरा.
\end{editorial}

\olimport{rules-and-proofs}

\olimport{propositional-rules}

\iftag{FOL}{%
  \olimport{quantifier-rules}
}{}

\olimport{structural-rules}

\olimport{derivations}

\olimport{proving-things}

\iftag{FOL}{%
  \olimport{proving-things-quant}
}{}

\olimport{proof-theoretic-notions}

\olimport{provability-consistency}

\olimport{provability-propositional}

\iftag{FOL}{%
  \olimport{provability-quantifiers}
}{}

\olimport{soundness}

\iftag{FOL}{%
  \olimport{identity}
  \olimport{soundness-identity}
}{}

\OLEndChapterHook

\end{document}
